\section{Introduction}
\label{sec:introduction}

Large language model (LLM) agents are moving from single-session assistants to persistent systems that act across long horizons, multiple tasks, and evolving user states.
In such settings, future decisions depend on past interaction: user preferences, task commitments, tool outcomes, failed attempts, environmental changes, and previously learned procedures.
Scaling the context window is an incomplete solution.
Even when longer histories can be provided, relevant information may be buried, outdated, displaced by irrelevant context, or present but unused.
Long-horizon agents therefore need mechanisms that make past interaction usable at decision time.

Recent work has made substantial progress on this problem by externalizing interaction history into reusable substrates: user states, episodic records, relational stores, graphs, semantic workspaces, skills, rules, and compressed experience artifacts.
Experience-compression frameworks make this trend explicit: raw interaction traces can be compressed into artifacts at different abstraction levels, trading specificity for reusability and lower context cost \cite{zhang2026experiencecompressionspectrumunifying}.
Runtime access has also moved beyond flat top-$k$ retrieval.
Modern systems combine retrieval with reranking, graph traversal, hierarchical selection, source reading, tool execution, workspace navigation, and trace-based attribution.
These mechanisms improve the geometry of access: they make relevant artifacts easier to locate, select, and expose.

These developments sharpen the access layer, but they also expose its limit.
The field has become good at making evidence easier to find.
It has not solved what happens after evidence is found.
In long-term personalized memory QA, aggregation-heavy questions can fail under oracle-evidence conditions, where the correct evidence is already present.
This is not a retrieval failure.
It is a post-access failure: existing systems optimize the geometry of access, while hard memory use also requires the algebra of evaluation.

We study this post-access regime.
Access determines which evidence reaches the model; post-access evaluation determines whether that evidence is exposed in a decision-sufficient view, whether variables are stably bound across evidence items, and whether the final answer is constrained to the evidence-defined valid set.
A flattened textual view may present a timestamp and a caption while omitting the relations, quantities, source pointers, or modality-specific structure needed for aggregation.
The evidence is in context, but the view is not decision-sufficient.

This distinction matters because language-only decoding is a weak evaluator over recalled evidence.
Without explicit variable binding and constraint enforcement, the model can match, paraphrase, and classify fluently, but it does not reliably execute multi-evidence aggregation, conflict resolution, or discrete computation.
Soft controls such as instruction tuning, prompting, chain-of-thought, and schema formatting can bias the model toward aggregation-like trajectories, but they do not enforce the answer set.

We use Artifact Access Semantics to organize the runtime path from compressed artifacts to access operations, primitive traces, evidence views, and decisions.
Within this path, our focus is the post-access boundary where recalled evidence must become an answerable state.
We formalize the access--evaluation separation and instantiate it in long-term personalized memory QA under controlled evidence conditions.
Our thesis is simple: multimodal long-term memory systems need both evidence access and evidence aggregation.
Existing systems have advanced the first substantially; the second---view construction, variable binding, operator execution, and evidence-valid decoding---remains the bottleneck.
